\section{模式三：完整轮次评测 --- 端到端验证}

\subsection{完整轮次的三个检查维度}

完整轮次评测让 Agent 从接收输入到返回最终结果完整运行一次。与单步评测的``单元测试''类比不同，完整轮次更像是\textbf{集成测试}，关注 Agent 多步协作后的整体效果。LangChain 总结了三个检查维度：

\subsubsection{轨迹检查（Trajectory）}

验证 Agent 在整个执行过程中是否调用了必要的工具，\textbf{但不要求严格的调用顺序}。例如日历调度 Agent 可能需要多次工具调用才能找到所有参与者都空闲的时间段——我们只关心它最终是否查询了所有相关日历，而不关心查询的顺序。

\subsubsection{最终响应检查（Final Response）}

对于编程和研究等开放性任务，Agent 到达目标的\textbf{路径}可能有很多种，此时最终输出的质量比具体路径更重要。这与传统的``标准答案''评测模式有根本区别。

\subsubsection{其他状态检查（Other State）}

很多 Deep Agent 不仅仅产生文本回复，还会创建 artifacts（文件、代码、配置等）。这些副产物也需要纳入评测范围：

\begin{table}[H]
\centering
\caption{不同类型 Agent 的状态检查策略}
\begin{tabular}{lp{9cm}}
\toprule
\textbf{Agent 类型} & \textbf{状态检查方式} \\
\midrule
编程 Agent & 读取并测试 Agent 生成的代码文件（运行单元测试、检查语法等） \\
研究 Agent & 验证 Agent 是否找到了正确的链接或信息来源 \\
邮件 Agent & 检查 Agent 是否正确更新了用户偏好文件 \\
调度 Agent & 验证日历中是否实际创建了正确的事件 \\
\bottomrule
\end{tabular}
\end{table}

\begin{warningbox}{轨迹检查不要过度约束}
常见错误是对工具调用的\textbf{顺序}做严格断言。Agent 可能因为模型版本更新、temperature 差异等原因改变工具调用顺序，但只要最终结果正确，这种顺序变化是可以接受的。过度约束会导致\textbf{脆弱测试（flaky tests）}，增加维护成本。建议使用 \texttt{any()} 而非 \texttt{assertEqual(trajectory, expected\_trajectory)} 来做轨迹断言。
\end{warningbox}

\subsection{LangSmith Trace 可视化}

LangSmith 提供了完整轮次的 trace 可视化能力，可以查看：

\begin{itemize}
  \item 高层指标：延迟（latency）和 token 消耗
  \item 细粒度步骤：每次模型调用和工具调用的详细信息
  \item 失败定位：当测试失败时，trace 视图可以快速定位到具体出错的步骤
\end{itemize}

这种从宏观到微观的多层级追踪对于调试复杂 Agent 行为至关重要。

\subsection{本章小结}

完整轮次评测提供了 Agent 行为的全景视图。通过轨迹、最终响应和其他状态三个维度的检查，可以全面评估 Agent 的端到端表现。关键原则是：\textbf{检查结果而非路径}，对轨迹做适度宽松的断言以避免脆弱测试。


